\documentclass[12pt,a4paper]{article}

\usepackage[utf8]{inputenc}
\usepackage[T1]{fontenc}
\usepackage{lmodern}
\usepackage{geometry}
\geometry{margin=1in}
\usepackage{graphicx}
\usepackage{amsmath, amssymb}
\usepackage{hyperref}
\usepackage{cancel}
\usepackage{color}

\title{DCON $recon flag$ document}
\author{Sunjae Lee}
\date{\today}

\begin{document}

\maketitle

% \begin{abstract}
% This is the abstract of the document. Briefly describe the purpose and content here.
% \end{abstract}

\section{Definition of covariant and contravariant basis}
The dcon context will be in black and \color{blue}the reference\cite{chance1992mhd} would be in blue. 
\color{black}\subsection{Covariant basis}
\begin{equation}
    \begin{aligned}
    \vec{v}_1 &= \vec{\nabla}\theta \times \vec{\nabla}\zeta, &
    \vec{v}_2 &= \vec{\nabla}\zeta \times \vec{\nabla}\psi, &
    \vec{v}_3 &= \vec{\nabla}\psi \times \vec{\nabla}\theta, 
    \end{aligned}
\end{equation}

\begin{equation}
    \begin{aligned}
    \vec{v}_i &= v_{ij}\hat{e}_j, \qquad &
    % g_{ij} &\equiv \color{red}\mathcal{J}\,\vec{v}_i \cdot \vec{v}_j = \mathcal{J}\,v_{ik}v_{kj}
    \end{aligned}
\end{equation}


\begin{equation}
\begin{aligned}
v_{11} &= \vec{v}_1 \cdot \hat{e}_1 = \frac{1}{2r\mathcal{J}}\frac{\partial r^2}{\partial \psi}, \\[6pt]
v_{12} &= \vec{v}_1 \cdot \hat{e}_2 = \frac{2\pi r}{\mathcal{J}}\frac{\partial \eta}{\partial \psi}, \\[6pt]
v_{13} &= \vec{v}_1 \cdot \hat{e}_3 = \frac{R}{\mathcal{J}}\frac{\partial \phi}{\partial \psi}, , \\[6pt]
\end{aligned}
\end{equation}

\begin{equation}
\begin{aligned}
v_{21} &= \vec{v}_2 \cdot \hat{e}_1 = \frac{1}{2r\mathcal{J}}\frac{\partial r^2}{\partial \theta}, \\[6pt]
v_{22} &= \vec{v}_2 \cdot \hat{e}_2 = \frac{2\pi r}{\mathcal{J}}\frac{\partial \eta}{\partial \theta}, \\[6pt]
v_{23} &= \vec{v}_2 \cdot \hat{e}_3 = \frac{R}{\mathcal{J}}\frac{\partial \phi}{\partial \theta}, \\[6pt]
\end{aligned}
\end{equation}

\begin{equation}
\begin{aligned}
v_{31} &= \vec{v}_3 \cdot \hat{e}_1 = \frac{1}{2r\mathcal{J}}\frac{\partial r^2}{\partial \zeta} = 0, \\[6pt]
v_{32} &= \vec{v}_3 \cdot \hat{e}_2 = \frac{2\pi r}{\mathcal{J}}\frac{\partial \eta}{\partial \zeta} = 0, \\[6pt]
v_{33} &= \vec{v}_3 \cdot \hat{e}_3 = \frac{R}{\mathcal{J}}\frac{\partial \phi}{\partial \zeta} = \frac{2\pi R}{\mathcal{J}}
\end{aligned}
\end{equation}


\subsubsection{Contravariant basis}
\begin{equation}
\begin{aligned}
\vec{w}_1 &= \vec{\nabla}\psi, &
\vec{w}_2 &= \vec{\nabla}\theta, &
\vec{w}_3 &= \vec{\nabla}\zeta, 
\end{aligned}
\end{equation}

\begin{equation}
    \begin{aligned}
    \vec{w}_i &= w_{ij}\hat{e}_j, \qquad &
    \mathcal{J}\,\vec{v}_i \cdot \vec{w}_j &= \delta_{ij}.
    \end{aligned}
\end{equation}

% ==== All w_{ij}: cross-form -> v-components ====
% Conventions:
%  v_1 = ∇θ × ∇ζ,  v_2 = ∇ζ × ∇ψ,  v_3 = ∇ψ × ∇θ
%  w_1 = ∇ψ,       w_2 = ∇θ,       w_3 = ∇ζ
%  v_i = v_{ij} \hat e_j,  w_i = w_{ij} \hat e_j
%  General identity:  w_{i\alpha} = J\, \varepsilon_{\alpha\beta\gamma}\, v_{j\beta} v_{k\gamma}
%  with (i,j,k) cyclic: (1,2,3), (2,3,1), (3,1,2)
\begin{equation}
    \begin{aligned}
    % i = 1
    w_{11}
    &= \vec\nabla\psi \cdot \hat e_1
    = \mathcal{J}\,\Big[\,(\vec\nabla\zeta \times \vec\nabla\psi)\times(\vec\nabla\psi \times \vec\nabla\theta)\,\Big]\!\cdot\hat e_1
    = \mathcal{J}\,(v_{22}v_{33}-v_{23}v_{32})
    = \frac{4\pi^{2} r R}{\mathcal{J}}\frac{\partial \eta}{\partial \theta}, \\[6pt]
    w_{12}
    &= \vec\nabla\psi \cdot \hat e_2
    = \mathcal{J}\,\Big[\,(\vec\nabla\zeta \times \vec\nabla\psi)\times(\vec\nabla\psi \times \vec\nabla\theta)\,\Big]\!\cdot\hat e_2
    = \mathcal{J}\,(v_{23}v_{31}-v_{21}v_{33})
    = -\,\frac{\pi R}{r\,\mathcal{J}}\frac{\partial r^{2}}{\partial \theta}, \\[6pt]
    w_{13}
    &= \vec\nabla\psi \cdot \hat e_3
    = \mathcal{J}\,\Big[\,(\vec\nabla\zeta \times \vec\nabla\psi)\times(\vec\nabla\psi \times \vec\nabla\theta)\,\Big]\!\cdot\hat e_3
    = \mathcal{J}\,(v_{21}v_{32}-v_{22}v_{31})
    = 0
    \end{aligned}
\end{equation}

\begin{equation}
    \begin{aligned}
    w_{21}
    &= \vec\nabla\theta \cdot \hat e_1
    = \mathcal{J}\,\Big[\,(\vec\nabla\psi \times \vec\nabla\theta)\times(\vec\nabla\theta \times \vec\nabla\zeta)\,\Big]\!\cdot\hat e_1
    = \mathcal{J}\,(v_{32}v_{13}-v_{33}v_{12})
    = -\,\frac{4\pi^{2} r R}{\mathcal{J}}\frac{\partial \eta}{\partial \psi}, \\[6pt]
    w_{22}
    &= \vec\nabla\theta \cdot \hat e_2
    = \mathcal{J}\,\Big[\,(\vec\nabla\psi \times \vec\nabla\theta)\times(\vec\nabla\theta \times \vec\nabla\zeta)\,\Big]\!\cdot\hat e_2
    = \mathcal{J}\,(v_{33}v_{11}-v_{31}v_{13})
    = \frac{\pi R}{r\,\mathcal{J}}\frac{\partial r^{2}}{\partial \psi}, \\[6pt]
    w_{23}
    &= \vec\nabla\theta \cdot \hat e_3
    = \mathcal{J}\,\Big[\,(\vec\nabla\psi \times \vec\nabla\theta)\times(\vec\nabla\theta \times \vec\nabla\zeta)\,\Big]\!\cdot\hat e_3
    = \mathcal{J}\,(v_{31}v_{12}-v_{32}v_{11})
    = 0
    \end{aligned}
\end{equation}

\begin{equation}
    \begin{aligned}
    w_{31}
    &= \vec\nabla\zeta \cdot \hat e_1
    = \mathcal{J}\,\Big[\,(\vec\nabla\theta \times \vec\nabla\zeta)\times(\vec\nabla\zeta \times \vec\nabla\psi)\,\Big]\!\cdot\hat e_1
    = \mathcal{J}\,(v_{12}v_{23}-v_{13}v_{22})
    = \frac{2\pi r R}{\mathcal{J}}\!\left(
        \frac{\partial \eta}{\partial \psi}\frac{\partial \phi}{\partial \theta}
    - \frac{\partial \phi}{\partial \psi}\frac{\partial \eta}{\partial \theta}
    \right), \\[6pt]
    w_{32}
    &= \vec\nabla\zeta \cdot \hat e_2
    = \mathcal{J}\,\Big[\,(\vec\nabla\theta \times \vec\nabla\zeta)\times(\vec\nabla\zeta \times \vec\nabla\psi)\,\Big]\!\cdot\hat e_2
    = \mathcal{J}\,(v_{13}v_{21}-v_{11}v_{23})
    = \frac{R}{2r\,\mathcal{J}}\!\left(
        \frac{\partial \phi}{\partial \psi}\frac{\partial r^{2}}{\partial \theta}
    - \frac{\partial r^{2}}{\partial \psi}\frac{\partial \phi}{\partial \theta}
    \right), \\[6pt]
    w_{33}
    &= \vec\nabla\zeta \cdot \hat e_3
    = \mathcal{J}\,\Big[\,(\vec\nabla\theta \times \vec\nabla\zeta)\times(\vec\nabla\zeta \times \vec\nabla\psi)\,\Big]\!\cdot\hat e_3
    = \mathcal{J}\,(v_{11}v_{22}-v_{12}v_{21})
    = \frac{1}{2\pi R}.
    \end{aligned}
\end{equation}

\section{Coordinate system}

The coordinate system comparing between DCON and reference\cite{chance1992mhd}.
In reference \cite{chance1992mhd}, $2 \pi \psi$ is poloidal flux 
\begin{equation}
\Psi_p(\psi) = 2\pi \psi 
= \int_{0}^{2\pi} d\phi 
  \int_{\text{axis}}^{\psi} 
  \mathbf{B}_p \cdot d\mathbf{S}
\end{equation}

In DCON, $\psi_{dcon}$, $\theta_{dcon}$, $\zeta_{dcon}$ are all normalized in [0,1]. $\chi_{dcon}$ representing the poloidal flux function and $\alpha_{dcon}$ for SFL angle. Also to scale the actual poloidal flux $\chi\prime_{dcon}$ is used.($\chi\prime_{dcon} = 2\pi(ssibry-ssimag)$)

\begin{align}
\mathbf{B} &= \nabla \chi_{dcon} \times \nabla \alpha_{dcon}\\
%
&= \chi\prime_{dcon} \nabla \psi_{dcon} \times \nabla \alpha_{dcon}\\
%
&= \chi\prime_{dcon} \nabla \psi_{dcon} \times (q\nabla \theta - \nabla \zeta)\\
%
&= f_{dcon} \nabla \phi + \nabla \phi \times \frac{\chi\prime_{dcon}}{2\pi} \nabla \psi_{dcon} \\
%
&= \color{blue} \nabla \psi \times (- \nabla \alpha)\\
%
&= \color{blue} g(\psi) \nabla \phi + \nabla \phi \times \nabla \psi \\
%
\end{align}

\begin{equation}
    \alpha_{dcon} = \color{blue} {- \alpha / 2 \pi} \color{black}= q(\psi)\theta - \zeta 
\end{equation}
\begin{equation}
g(\psi) \leftrightarrow (toroidal field functions) \leftrightarrow f_{dcon}
\end{equation}

So as follows, \[
\color{blue}{\psi}
\;\;\longleftrightarrow\;\;
\color{black}\frac{\chi_{\mathrm{dcon}}}{2\pi}
\]
%
\[
\color{blue}{\psi}
= \color{black}\frac{\chi'_{\mathrm{dcon}}}{2\pi}\,
\psi_{\mathrm{dcon}}.
\]
will established.
And also 
\[
{\chi} = \color{black}\chi'_{\mathrm{dcon}} \psi_{\mathrm{dcon}} 
\]

\begin{align}
\nabla \cdot \left( \frac{1}{R^2} \nabla \psi \right) 
&= \nabla \cdot \left(\frac{1}{R^2}\, \frac{\chi\prime_{dcon}}{2\pi} \nabla\, \psi_{dcon} \right)\\
&= \frac{1}{{2\pi}} \nabla \cdot \left(\frac{1}{R^2}\, \chi\prime_{dcon} \nabla\, \psi_{dcon} \right) \\
&= \frac{1}{{2\pi}} \nabla \cdot \left(\frac{1}{R^2}\, \nabla \chi_{dcon} \right) \\
&= \mathbf{J} \cdot \nabla \phi\\
&= - \left( p'_{chance} + \frac{1}{R^2} g g' \right)
\end{align}

\begin{equation}
    R^2 \nabla \cdot \left( \frac{1}{R^2} \nabla \chi_{dcon} \right) = - 2\pi \left(R^2 p'_{chance} + g g' \right) = -(2\pi)\frac{2\pi}{\chi\prime_{dcon}}(ff'+ R^2p')
\end{equation}



\subsection{Derivation of Eq.~(29) in GPEC note}

\begin{equation}
    \begin{aligned}
    B 
    &= \chi\prime \nabla \psi \times (q\nabla \theta - \nabla \zeta) = \chi\prime q \vec{v}_3 + \chi\prime \vec{v}_2 \\
    &= \mathcal{J}(B^\psi \vec{v}_1 + B^\theta \vec{v}_2 + B^\zeta \vec{v}_3) \\
    &= (B_\psi \vec{w}_1 + B_\theta \vec{w}_2 + B_\zeta \vec{w}_3) 
\end{aligned}
\end{equation}

% \begin{equation}
% S \;=\; -\,\frac{\vec{B} \times \vec{\nabla}\psi}{\lvert \vec{\nabla}\psi \rvert^{2}}
% \;\cdot\;
% \left(\vec{\nabla} \times 
% \frac{\vec{B} \times \vec{\nabla} \psi}{\lvert \vec{\nabla} \psi \rvert^{2}}\right)
% \end{equation}

\begin{equation}
\begin{aligned}
B^\theta &\equiv \vec{B}\cdot \vec{w}_2 = \frac{\chi'}{\mathcal{J}},
\qquad
B^\zeta \equiv \vec{B}\cdot \vec{w}_3 = \frac{q\,\chi'}{\mathcal{J}}, \\[4pt]
\vec{B} \times \vec{w}_\psi
&= \frac{1}{\mathcal{J}}\Big( B^\zeta\,\vec{v}_2 - B^\theta\,\vec{v}_3 \Big)
= \frac{1}{\mathcal{J}}\left(\frac{q\,\chi'}{\mathcal{J}}\,\vec{v}_2 - \frac{\chi'}{\mathcal{J}}\,\vec{v}_3\right) \\
&= \frac{\chi'}{\mathcal{J}^{2}}\Big( q\,\vec{v}_2 - \vec{v}_3 \Big),
\end{aligned}
\end{equation}

\begin{equation}
(\vec{B} \times \vec{w}_1)^\psi = 0, \qquad
(\vec{B} \times \vec{w}_1)^\theta = \frac{q\,\chi'}{\mathcal{J}^{2}}, \qquad
(\vec{B} \times \vec{w}_1)^\zeta = -\,\frac{\chi'}{\mathcal{J}^{2}}.
\end{equation}

\subsection{Metric Tensor and Index Operations}

Given the definitions of the basis vectors where $\vec{v}_i = \vec{\nabla}\theta \times \vec{\nabla}\zeta$ (requiring the Jacobian $\mathcal{J}$ for physical dimensions) and $\vec{w}_i = \vec{\nabla}\psi$ (representing the dual basis directly), the metric tensors for lowering and raising indices are derived as follows.

\subsubsection{Lowering Operator ($B^j \to B_i$)}
To transform contravariant components $B^j$ into covariant components $B_i$, we utilize the relation between the physical basis $\vec{e}_i$ and the defined basis $\vec{v}_i$, where $\vec{e}_i = \mathcal{J}\vec{v}_i$.

Using the definition $\vec{B} = \mathcal{J} B^j \vec{v}_j$ and taking the dot product with $\vec{e}_i$:
\begin{equation}
    \begin{aligned}
    B_i &= \vec{B} \cdot \vec{e}_i \\
        &= (\mathcal{J} B^j \vec{v}_j) \cdot (\mathcal{J} \vec{v}_i) \\
        &= \mathcal{J}^2 (\vec{v}_i \cdot \vec{v}_j) B^j.
    \end{aligned}
\end{equation}
Thus, the covariant metric tensor $g_{ij}$ must include the square of the Jacobian:
\begin{equation}
    g_{ij} \equiv \mathcal{J}^2 (\vec{v}_i \cdot \vec{v}_j).
\end{equation}

\subsubsection{Raising Operator ($B_j \to B^i$)}
To transform covariant components $B_j$ into contravariant components $B^i$, we use the gradient basis $\vec{w}_i$, which corresponds to the dual basis $\vec{e}^i$ (i.e., $\vec{e}^i = \vec{w}_i$).

Using the definition $\vec{B} = B_j \vec{w}_j$ and taking the dot product with $\vec{w}_i$:
\begin{equation}
    \begin{aligned}
    B^i &= \vec{B} \cdot \vec{w}_i \\
        &= (B_j \vec{w}_j) \cdot \vec{w}_i \\
        &= (\vec{w}_i \cdot \vec{w}_j) B_j.
    \end{aligned}
\end{equation}
Unlike the lowering operator, the contravariant metric tensor $g^{ij}$ does not require the Jacobian:
\begin{equation}
    g^{ij} \equiv \vec{w}_i \cdot \vec{w}_j.
\end{equation}

\subsection{Derivation of Eq. (35) in gpec note}


\[
w_{11} = \frac{4 \pi^2 r R}{\mathcal{J}} \frac{\partial \eta}{\partial \theta},
\quad
v_{22} = \frac{2 \pi r}{\mathcal{J}} \frac{\partial \eta}{\partial \theta},
\quad
w_{12} = \frac{-\pi R}{r \mathcal{J}} \frac{\partial r^2}{\partial \theta},
\quad
v_{21} = \frac{1}{2r \mathcal{J}} \frac{\partial r^2}{\partial \theta}
\]

\[
w_{11} = 2 \pi R v_{22}, 
\qquad
w_{12} = -2 \pi R v_{21}
\]


\[
q\bigl({\vec{\nabla}} \psi \cdot {\vec{\nabla}} \theta \bigr)
- {\vec{\nabla}} \psi \cdot {\vec{\nabla}} \zeta 
\]
\[
= q\,(w_{11} w_{21} + w_{12} w_{22})
 - (w_{11} w_{31} + w_{12} w_{32})
\]

\[
= w_{11} (q w_{21} - w_{31}) + w_{12} (q w_{22} - w_{32})
\]

\[
= 2\pi R \Bigl[
    q v_{22} w_{21}
  - q v_{21} w_{22}
  - v_{22} w_{31}
  + v_{21} w_{32}
\Bigr]
\]

\[
= 2\pi R \mathcal{J} \Bigl[
    q v_{22} (\cancel{v_{32}v_{13}} - v_{33} v_{12})
  - q v_{21} (v_{33} v_{11} - \cancel{v_{31}v_{13}})
  - v_{22}(v_{23} v_{12} - v_{13} v_{22})
  + v_{21}(v_{13} v_{21} - v_{11} v_{23})
\Bigr]
\]


\[
= 2\pi R \mathcal{J} \Bigl[
   - q v_{33} (v_{22} v_{12} + v_{21} v_{11})
   + v_{13}(v_{22}^2 + v_{21}^2)
   - v_{23}(v_{22} v_{12} + v_{11} v_{21})
\Bigr]
\]

\[
= 2\pi R \mathcal{J} \Bigl[
   - (v_{22} v_{12} + v_{21} v_{11})(v_{23}+q v_{33})
   + v_{13}(v_{22}^2 + v_{21}^2)
\Bigr]
\]

\section{Energy principle\cite{RevModPhys.76.1071}\cite{Park2009Ideal}} 

\begin{align}
\delta W_p &= \frac{1}{2} \int d\tau_p 
\Big\{ |\mathbf{Q}|^2 - \mathbf{j} \cdot \, \mathbf{Q} \times \boldsymbol{\xi}^* 
+ \gamma P |\nabla \cdot \boldsymbol{\xi}|^2 
+ (\nabla \cdot \boldsymbol{\xi})^* \boldsymbol{\xi} \cdot \nabla P \Big\} \\
&= \frac{1}{2} \int d\tau_p 
\Big\{ |\mathbf{Q}_\perp|^2 
+ \frac{1}{B^2} |\mathbf{Q} \cdot \mathbf{B} - \boldsymbol{\xi} \cdot \nabla P|^2
+ \sigma\, \mathbf{B} \times \boldsymbol{\xi}^* \cdot \mathbf{Q}
+ \gamma P |\nabla \cdot \boldsymbol{\xi}|^2 
- 2 \boldsymbol{\xi} \cdot \nabla P\, \boldsymbol{\xi}^* \cdot \boldsymbol{\kappa}
\Big\} \\
&= \frac{1}{2} \int d\tau_p 
\Big\{ |\mathbf{Q} + \xi_n \mathbf{j} \times \mathbf{n}|^2 
+ \gamma P |\nabla \cdot \boldsymbol{\xi}|^2 
- 2\,\mathbf{j} \times \mathbf{n} \cdot (\mathbf{B} \cdot \nabla \mathbf{n})\, \xi_n^2
\Big\}\\
&= \frac{1}{2} \int d\tau_p 
\Big\{ |\mathbf{Q} + \xi_n j \times \mathbf{n}|^2 
+ \gamma P |\nabla \cdot \boldsymbol{\xi}|^2 
- K\, \xi_n^2
\Big\}
\end{align}

The derivation of (35) could be found in \cite{RevModPhys.76.1071} and the original derivation of (36) is in \cite{Bernstein1983}.
The K term in (37) is defined in \cite{chance1992mhd}. 


\begin{equation}
\sigma = \frac{\mathbf{j} \cdot \mathbf{B}}{B^2}, \hat{\mathbf n}=\frac{\nabla\psi}{|\nabla\psi|}, \mathbf Q \equiv \nabla\times(\boldsymbol{\xi}\times\mathbf B) (\delta\mathbf B = \mathbf Q)
\end{equation}

\subsection{Transformation of $\delta W_F$ in \cite{Bernstein1983}(Modified units)}
In this section, we will see how (34) changes to (36). If we define $\xi$ as follows:
\begin{equation}
\boldsymbol{\xi}=a\,\nabla p + b\,\nabla\times\mathbf{B} + c\,\mathbf{B}
\end{equation}
then
\begin{equation}
\boldsymbol{\xi}\times\mathbf{B}
= -a\,\mathbf{B}\times\nabla p + \mu_0 b\,\nabla p
\end{equation}
\begin{equation}
\boldsymbol{\xi}\times(\nabla\times\mathbf{B})
= -a\,(\nabla\times\mathbf{B})\times\nabla p - \mu_0 c\,\nabla p
\end{equation}

Using the SI form of $\delta W_F$ ($1/4\pi\to 1/\mu_0$),
Gauss's theorem, and the boundary condition $\mathbf{n}\cdot\mathbf{B}=0$,
one finds
\begin{align}
2\mu_0 W_F
& = \int d^3r\Bigl\{ \mu_0\gamma p(\nabla\cdot\boldsymbol{\xi})^2 
+ |\mathbf{Q}|^2 \nonumber + (\nabla \cdot \boldsymbol{\xi})^* \boldsymbol{\xi} \cdot \nabla P  \nonumber
+ \boldsymbol{\xi} \times \mathbf{Q} \cdot \nabla \times \mathbf{B}
\Bigr\}\\
& = \int d^3r\Bigl\{\mu_0\gamma p(\nabla\cdot\boldsymbol{\xi})^2
+ \bigl[\nabla\times\bigl(-a\,\mathbf{B}\times\nabla p + \mu_0 b\,\nabla p\bigr)\bigr]^2
\nonumber\\
&+ \mu_0 a(\nabla p)^2\nabla\cdot\bigl(a\nabla p + b\nabla\times\mathbf{B}+c\mathbf{B}\bigr)
\nonumber\\
&+ \bigl[a(\nabla\times\mathbf{B})\times\nabla p + \mu_0 c\,\nabla p\bigr]\cdot
\nabla\times\bigl(-a\,\mathbf{B}\times\nabla p + \mu_0 b\,\nabla p\bigr)
\Bigr\}.
\end{align}
Since,
\begin{align}
&a(\nabla p)^2\nabla\cdot\bigl(b\nabla\times\mathbf{B}+c\mathbf{B}\bigr)
-\nabla\cdot\Bigl[a(\nabla p)^2\bigl(b\nabla\times\mathbf{B}+c\mathbf{B}\bigr)\Bigr]
\nonumber\\
&= -\bigl(b\nabla\times\mathbf{B}+c\mathbf{B}\bigr)\cdot\nabla\!\Bigl[a(\nabla p)^2\Bigr]
\nonumber\\
&= -b\,\nabla\cdot\Bigl[a(\nabla p)^2(\nabla\times\mathbf{B})\Bigr]
   -c\,\nabla\cdot\Bigl[a(\nabla p)^2\mathbf{B}\Bigr]
\nonumber\\
&= -b\,\nabla\cdot\Bigl[(\nabla p)\times\bigl(a(\nabla\times\mathbf{B})\times\nabla p\bigr)\Bigr]
   -c\,\nabla\cdot\Bigl[(\nabla p)\times\bigl(a\mathbf{B}\times\nabla p\bigr)\Bigr]
\nonumber\\
&= b(\nabla p)\cdot\nabla\times\Bigl[a(\nabla\times\mathbf{B})\times\nabla p\Bigr]
 + c(\nabla p)\cdot\nabla\times\Bigl[a\mathbf{B}\times\nabla p\Bigr]
\nonumber\\
&= \nabla\cdot\Bigl(\bigl[a(\nabla\times\mathbf{B})\times\nabla p\bigr]\times b\nabla p\Bigr)
 + a(\nabla\times\mathbf{B})\times(\nabla p)\cdot\nabla\times(b\nabla p)
\nonumber\\
&\qquad + c(\nabla p)\cdot\nabla\times(a\mathbf{B}\times\nabla p).
\end{align}
so,
\begin{align}
a(\nabla p)^2\nabla\cdot\bigl(b\nabla\times\mathbf{B}+c\mathbf{B}\bigr)
&= \nabla\cdot\Bigl[ac(\nabla p)^2\mathbf{B}\Bigr]
 + a(\nabla\times\mathbf{B})\times(\nabla p)\cdot\nabla\times(b\nabla p)
\nonumber\\
&\quad + c(\nabla p)\cdot\nabla\times(a\mathbf{B}\times\nabla p).
\end{align}

Using (44) to eliminate the divergence term involving $b,c$ gives
\begin{align}
2\mu_0 W_F
= \int d^3r\Bigl\{
&\mu_0\gamma p(\nabla\cdot\boldsymbol{\xi})^2
+ \bigl[\nabla\times\bigl(-a\,\mathbf{B}\times\nabla p + \mu_0 b\,\nabla p\bigr)\bigr]^2
\nonumber\\
&+ \mu_0 a(\nabla p)^2\nabla\cdot(a\nabla p)
+ 2\mu_0 a(\nabla\times\mathbf{B})\times(\nabla p)\cdot\nabla\times(b\nabla p)
\nonumber\\
&- a(\nabla\times\mathbf{B})\times(\nabla p)\cdot\nabla\times(a\mathbf{B}\times\nabla p)
\Bigr\}
\nonumber\\
= \int d^3r\Bigl\{
&\mu_0\gamma p(\nabla\cdot\boldsymbol{\xi})^2
+ \bigl[\nabla\times\bigl(-a\,\mathbf{B}\times\nabla p + \mu_0 b\,\nabla p\bigr)
      + a(\nabla\times\mathbf{B})\times\nabla p\bigr]^2
\nonumber\\
&+ a(\nabla\times\mathbf{B})\times(\nabla p)\cdot\nabla\times(a\mathbf{B}\times\nabla p)
- a^2\bigl[(\nabla\times\mathbf{B})\times\nabla p\bigr]^2
+ \mu_0 a(\nabla p)^2\nabla\cdot(a\nabla p)
\Bigr\}.
\end{align}

Define
\begin{equation}
\mathbf{n}=\frac{\nabla p}{|\nabla p|}.
\end{equation}
Then, since $\mu_0\nabla p=(\nabla\times\mathbf{B})\times\mathbf{B}$,
\begin{align}
&-a^2\bigl[(\nabla\times\mathbf{B})\times\nabla p\bigr]^2
+ a(\nabla\times\mathbf{B})\times(\nabla p)\cdot\nabla\times(a\mathbf{B}\times\nabla p)
+ \mu_0 a(\nabla p)^2\nabla\cdot(a\nabla p)
\nonumber\\
&= -a^2(\nabla\times\mathbf{B})^2(\nabla p)^2
+ a^2(\nabla\times\mathbf{B})\times(\nabla p)\cdot\nabla\times(\mathbf{B}\times\nabla p)
\nonumber\\
&\quad + a(\nabla\times\mathbf{B})\times(\nabla p)\cdot(\nabla a)\times(\mathbf{B}\times\nabla p)
+ \mu_0 a^2(\nabla p)^2\nabla^2 p
+ \mu_0 a(\nabla p)^2(\nabla a)\cdot\nabla p
\nonumber\\
&= -a^2(\nabla\times\mathbf{B})^2(\nabla p)^2
+ a^2(\nabla\times\mathbf{B})\times(\nabla p)\cdot\Bigl[(\nabla p)\cdot\nabla\mathbf{B}
- \mathbf{B}\cdot\nabla\nabla p + \mathbf{B}\nabla^2 p\Bigr]
\nonumber\\
&\quad + a(\nabla\times\mathbf{B})\times(\nabla p)\cdot\Bigl[\mathbf{B}(\nabla a)\cdot\nabla p
- (\nabla p)\,\mathbf{B}\cdot\nabla a\Bigr]
+ \mu_0 a^2(\nabla p)^2\nabla^2 p
+ \mu_0 a(\nabla p)^2(\nabla a)\cdot\nabla p
\nonumber\\
&= a^2(\nabla p)^2\Bigl[-(\nabla\times\mathbf{B})^2
+ (\nabla\times\mathbf{B})\times\mathbf{n}\cdot(\mathbf{n}\cdot\nabla\mathbf{B}
- \mathbf{B}\cdot\nabla\mathbf{n})\Bigr].
\end{align}

Since $\nabla\cdot\mathbf{B}=0$, it follows that
\begin{equation}
0=\nabla(\mathbf{n}\cdot\mathbf{B})
=\mathbf{n}\cdot\nabla\mathbf{B}+\mathbf{B}\cdot\nabla\mathbf{n}
+\mathbf{n}\times(\nabla\times\mathbf{B})+\mathbf{B}\times(\nabla\times\mathbf{n}),
\end{equation}
whence
\begin{align}
&-(\nabla\times\mathbf{B})^2
+(\nabla\times\mathbf{B})\times\mathbf{n}\cdot\bigl[\mathbf{n}\cdot\nabla\mathbf{B}
-\mathbf{B}\cdot\nabla\mathbf{n}\bigr]
\nonumber\\
&=-(\nabla\times\mathbf{B})^2
+(\nabla\times\mathbf{B})\times\mathbf{n}\cdot\bigl[-2\mathbf{B}\cdot\nabla\mathbf{n}
-\mathbf{n}\times(\nabla\times\mathbf{B})-\mathbf{B}\times(\nabla\times\mathbf{n})\bigr]
\nonumber\\
&=-(\nabla\times\mathbf{B})^2+\bigl[\mathbf{n}\times(\nabla\times\mathbf{B})\bigr]^2
-2(\nabla\times\mathbf{B})\times\mathbf{n}\cdot(\mathbf{B}\cdot\nabla\mathbf{n})
\nonumber\\
&\quad -(\nabla\times\mathbf{B})\cdot\mathbf{B}\,\mathbf{n}\cdot\nabla\times\mathbf{n}
-(\nabla\times\mathbf{B})\cdot(\nabla\times\mathbf{n})\,\mathbf{n}\cdot\mathbf{B}
\nonumber\\
&=-2(\nabla\times\mathbf{B})\times\mathbf{n}\cdot(\mathbf{B}\cdot\nabla\mathbf{n}).
\end{align}

Since $\mu_0\nabla p=(\nabla\times\mathbf{B})\times\mathbf{B}$, it follows that
$\mathbf{n}\cdot\nabla\times\mathbf{B}=0$, whence
\begin{equation}
\bigl[\mathbf{n}\times(\nabla\times\mathbf{B})\bigr]^2
=\mathbf{n}^2(\nabla\times\mathbf{B})^2-(\mathbf{n}\cdot\nabla\times\mathbf{B})^2
=(\nabla\times\mathbf{B})^2.
\end{equation}

Also, since $\mathbf{n}$ is a unit normal to the surfaces $p=\text{const}$, for any area
contained in such a surface, Stokes' theorem gives
\begin{equation}
\int d^2r\,\mathbf{n}\cdot\nabla\times\mathbf{n}
=\oint d\mathbf{r}\cdot\mathbf{n}=0,
\end{equation}
and since the surface is arbitrary, $\mathbf{n}\cdot\nabla\times\mathbf{n}=0$.

Thus one can write,
\begin{equation}
2\mu_0 W_F
=\int_I d^3r\Bigl\{
\mu_0\gamma p(\nabla\cdot\boldsymbol{\xi})^2
+\bigl[\mathbf{Q}+(\nabla\times\mathbf{B})\times\mathbf{n}\,(\mathbf{n}\cdot\boldsymbol{\xi})\bigr]^2
-2(\mathbf{n}\cdot\boldsymbol{\xi})^2\,(\nabla\times\mathbf{B})\times\mathbf{n}\cdot(\mathbf{B}\cdot\nabla\mathbf{n})
\Bigr\}.
\end{equation}

% If desired, replace curl B by current density j:
% \nabla\times\mathbf{B} = \mu_0\mathbf{j}, so the equilibrium is \nabla p=\mathbf{j}\times\mathbf{B}.
% \begin{align}
%     K &= 2(\mathbf{j} \times \hat{n}) \cdot (\mathbf{B} \cdot \nabla) \hat{n} \\
%     &=\frac{2}{|\nabla \psi|^2}(\mathbf{j} \times \nabla \psi) \cdot (\mathbf{B} \cdot \nabla) \nabla \psi \\
%     &=-(\frac{\mathbf{j} \cdot \mathbf{B}}{B^2})\mathbf{B} \times \nabla \psi\, \cdot\, (\nabla \times \frac{\mathbf{B} \times \nabla \psi}{|\nabla \psi|^2}) + \mathbf{j} \cdot \mathbf{B} + 2p'|\nabla \psi|^2 \kappa \cdot \nabla \psi \\
%     &=\frac{2}{|\nabla \chi|^2}(\mathbf{j} \times \nabla \psi) \cdot (\mathbf{B} \cdot \nabla) \nabla \chi   
% \end{align}
\subsection{Derivation of K}
\color{blue}
\begin{align}
K 
&= 2(\mathbf{j} \times \hat{n}) \cdot (\mathbf{B} \cdot \nabla) \hat{n}\\
&= \frac{2}{|\nabla \psi|} (\mathbf{j} \times \hat{n}) \cdot (\mathbf{B} \cdot \nabla) \nabla \psi + \cancel{2(\mathbf{j} \times \hat{n}) \cdot [\nabla \psi (\mathbf{B} \cdot \nabla) (\frac{1}{|\nabla \psi|})]}\\
&= -\frac{2}{|\nabla \psi|^2} (\mathbf{j} \times \nabla \psi) \cdot \big[(\nabla \psi \cdot \nabla)\mathbf{B} + \nabla \psi \times \mathbf{j}\big]\\
&= -\frac{2}{|\nabla \psi|^2} \big[(\mathbf{j} \times \nabla \psi)(\nabla \psi \cdot \nabla)\mathbf{B} - (\mathbf{j} \times \nabla \psi)^2\big]\\
&= +2j^2 -\frac{2}{|\nabla \psi|^2} \big[(\mathbf{j} \times \nabla \psi)(\nabla \psi \cdot \nabla)\mathbf{B} \big]\\
&= 2\sigma^2 B^2 + \underbrace{2 \frac{(\nabla p)^2}{B^2}-\frac{2}{|\nabla \psi|^2} \big[(\mathbf{j} \times \nabla \psi)(\nabla \psi \cdot \nabla)\mathbf{B} \big]}_{\text{$K_1 + K_3$}}
\end{align}
\color{black}
If we decompose current density into perpendicular and parallel components to magnetic field, then the perpendicular and parallel contributions to K are given by

\color{blue}
\begin{align}
K_3
&= -\frac{2}{|\nabla \psi|^2} (\mathbf{j}_\perp \times \nabla \psi) \cdot \big[(\nabla \psi \cdot \nabla)\mathbf{B}\big] +2 \frac{(\nabla p)^2}{B^2}\\
&= -\frac{2}{|\nabla \psi|^2} ((\frac{p^\prime}{B^2} \mathbf{B}\times \nabla \psi) \times \nabla \psi) \cdot \big[(\nabla \psi \cdot \nabla)\mathbf{B}\big] + 2 \frac{(p^\prime \nabla \psi)^2}{B^2}\\
&= \frac{2}{|\nabla \psi|^2}\frac{p^\prime}{B^2} [\mathbf{B} |\nabla \psi|^2 - \cancel{\nabla \psi (\nabla \psi \cdot \mathbf{B})}] \cdot \big[(\nabla \psi \cdot \nabla)\mathbf{B}\big]  + 2 \frac{(p^\prime \nabla \psi)^2}{B^2}\\
&= \frac{2p^\prime}{B^2} \{ \mathbf{B} \cdot [(\nabla \psi \cdot \nabla) \mathbf{B}] +(\nabla p \cdot \nabla \psi)\} \\
&= \frac{2p^\prime}{B^2} \{ \mathbf{B} \cdot [(\nabla \psi \cdot \nabla) \mathbf{B}] +(\mathbf{j} \times \mathbf{B}) \cdot \nabla \psi\} \\
&= \frac{2p^\prime}{B^2} \{ \mathbf{B} \cdot [(\nabla \psi \cdot \nabla) \mathbf{B}] + \mathbf{B} \cdot (\nabla \psi \times \mathbf{j})\}\\
&= -\frac{2p^\prime}{B^2} \mathbf{B} \cdot (\mathbf{B} \cdot \nabla) (\nabla \psi) = \frac{2p^\prime}{B^2}(\mathbf{B} \cdot \nabla) \cdot \mathbf{B} \cdot \nabla \psi \\
&= \frac{2p^\prime}{B^2}((\mathbf{B} \cdot \nabla) \cdot \mathbf{B} \cdot \nabla \psi - \nabla \frac{B^2}{2}\cdot \nabla \psi) = +2 p^\prime [ (\frac{\mathbf{B}}{B} \cdot \nabla) \frac{\mathbf{B}}{B} ]\cdot \nabla \psi \\
&= +2 p^\prime \boldsymbol{\kappa} \cdot \nabla \psi
\end{align}

% \color{black}
% \begin{equation}
% \mathbf{B}=B \mathbf{b}
% \end{equation}

\color{blue}
\begin{align}
K_1
&= -\frac{2}{|\nabla \psi|^2} (\mathbf{j}_\parallel \times \nabla \psi) \cdot \big[(\nabla \psi \cdot \nabla)\mathbf{B}\big] +\frac{2}{|\nabla \psi|^2}(\mathbf{j} \times \nabla \psi)^2\\
&= -\frac{2}{|\nabla \psi|^2} (\frac{\mathbf{j} \cdot \mathbf{B}}{|\mathbf{B}|^2} \mathbf{B}\times \nabla \psi) \cdot \big[(\nabla \psi \cdot \nabla)\mathbf{B}\big] \\
&= -\sigma \frac{2}{|\nabla \psi|^2} (\mathbf{B} \times \nabla \psi) \cdot \big[(\nabla \psi \cdot \nabla)\mathbf{B}\big] 
\end{align}

\color{black}
Since we can use the vector identity,
\begin{equation}
    \nabla\times(\mathbf a\times\mathbf c)
= \mathbf a(\nabla\cdot\mathbf c)-\mathbf c(\nabla\cdot\mathbf a)
+(\mathbf c\cdot\nabla)\mathbf a-(\mathbf a\cdot\nabla)\mathbf c.
\end{equation}
\begin{equation}
\nabla \times (\mathbf{B} \times \nabla \psi) = \mathbf{B} (\nabla^2 \psi) - \cancel{\nabla \psi (\nabla \cdot \mathbf{B}) }+ (\nabla \psi \cdot \nabla) \mathbf{B} - (\mathbf{B} \cdot \nabla) (\nabla \psi)
\end{equation}

\begin{equation}
    \nabla(\mathbf a\cdot\mathbf b)
    =(\mathbf a\cdot\nabla)\mathbf b+(\mathbf b\cdot\nabla)\mathbf a+\mathbf a\times(\nabla\times\mathbf b)+\mathbf b\times(\nabla\times\mathbf a)
\end{equation}
\begin{equation}
\cancel{\nabla(\mathbf B\cdot\nabla\psi)}=(\mathbf B\cdot\nabla)\nabla\psi+(\nabla\psi\cdot\nabla)\mathbf B+\nabla\psi\times(\nabla\times\mathbf B)+\cancel{\mathbf B\times(\nabla\times\nabla\psi)}
\end{equation}
\begin{equation}
    (\mathbf{B} \cdot \nabla) (\nabla \psi) = - (\nabla \psi \cdot \nabla) \mathbf{B} - \nabla \psi \times \mathbf{j}
\end{equation}

\begin{equation}
\nabla \times (\mathbf{B} \times \nabla \psi) = \mathbf{B} (\nabla^2 \psi) +2 (\nabla \psi \cdot \nabla) \mathbf{B} + \nabla \psi \times \mathbf{j}
\end{equation}
\begin{equation}
(\mathbf{B} \times \nabla \psi) \cdot [\nabla \times (\mathbf{B} \times \nabla \psi)] = \cancel{(\mathbf{B} \times \nabla \psi) \cdot \mathbf{B} (\nabla^2 \psi)} +2 (\mathbf{B} \times \nabla \psi) \cdot (\nabla \psi \cdot \nabla) \mathbf{B} + (\mathbf{B} \times \nabla \psi) \cdot (\nabla \psi \times \mathbf{j})
\end{equation}
and rearranging the terms, we have
\begin{equation}
    (\mathbf{B} \times \nabla \psi) \cdot \nabla \times (\frac{\mathbf{B} \times \nabla \psi}{|\nabla \psi|^2}) 
    = \cancel{(\mathbf{B} \times \nabla \psi) \cdot \nabla (\frac{1}{|\nabla \psi|^2}) \times (\mathbf{B} \times \nabla \psi) }+ (\mathbf{B} \times \nabla \psi) \cdot \frac{1}{|\nabla \psi|^2} \nabla \times (\mathbf{B} \times \nabla \psi)
\end{equation}

\color{blue}
\begin{align}
K_1 
&= -\sigma \frac{1}{|\nabla \psi|^2} (\mathbf{B} \times \nabla \psi) \cdot [\nabla \times (\mathbf{B} \times \nabla \psi)] +\sigma \frac{1}{|\nabla \psi|^2}(\mathbf{B} \times \nabla \psi) \cdot (\nabla \psi \times \mathbf{j})\\
&= -\sigma \mathbf{B}\times \nabla \psi\cdot (\nabla \times \frac{\mathbf{B} \times \nabla \psi}{|\nabla \psi|^2}) + \sigma \frac{1}{|\nabla \psi|^2} [-|\nabla \psi|^2 (\mathbf{B}\cdot\mathbf{j})+\cancel{(\mathbf{B}\cdot \nabla \psi)} \cdot (\nabla \psi \cdot \mathbf{j})]\\
&= \sigma S |\nabla \psi|^2 - \sigma (\mathbf{B}\cdot\mathbf{j}) = \sigma S |\nabla \psi|^2 - \sigma^2 B^2
\end{align}

\color{blue} 
\begin{align}
    K
    % &= - \mathbf{j} \cdot \mathbf{B} \frac{\mathbf{B}\times \nabla \psi}{|\mathbf{B}|^2} \cdot (\nabla \times \frac{\mathbf{B} \times \nabla \psi}{|\nabla \psi|^2}) + \mathbf{j} \cdot \mathbf{B} + 2p'(\frac{\mathbf{B}}{B} \cdot \nabla)\frac{\mathbf{B}}{B}\cdot \nabla \psi \\
    % &= - |\nabla \psi|^2 \sigma \frac{\mathbf{B}\times \nabla \psi}{|\nabla\psi|^2} \cdot (\nabla \times \frac{\mathbf{B} \times \nabla \psi}{|\nabla \psi|^2}) + \mathbf{j} \cdot \mathbf{B} + 2p'\boldsymbol{\kappa} \cdot \nabla \psi \\
    &= |\nabla \psi|^2 \sigma S + \sigma^2 B^2+ 2p'\boldsymbol{\kappa} \cdot \nabla \psi
\end{align}

\color{black} 
\subsection{Normalization of K between Chance and DCON coordinates}
\color{blue} 
\begin{align}
K&=|\nabla \psi|^2 \sigma S + \sigma \mathbf{j} \cdot \mathbf{B} + 2p'\boldsymbol{\kappa} \cdot \nabla \psi 
\end{align}

\begin{align}
    S&=-\frac{\mathbf{B}\times \nabla \psi}{|\nabla\psi|^2} \cdot (\nabla \times \frac{\mathbf{B} \times \nabla \psi}{|\nabla \psi|^2})
\end{align}
\color{black}

From this definition, the following relationships hold:
\begin{align}
\color{blue} |\nabla \psi|^2
&= \left(\frac{\chi^\prime_{\mathrm{dcon}}}{2\pi}\right)^2 |\nabla \psi_{\mathrm{dcon}}|^2
\end{align}

Substituting these relations into the original expression of \(K_{\mathrm{Chance}}\), we obtain
\begin{align}
K
&= \left(\frac{\chi^\prime_{\mathrm{dcon}}}{2\pi}\right)^2
|\nabla \psi_{\mathrm{dcon}}|^2 \sigma S
+ \sigma \mathbf{j}\!\cdot\!\mathbf{B}
+ 2p'\!\left(\!\boldsymbol{\kappa}\cdot\frac{\chi^\prime_{\mathrm{dcon}}}{2\pi}\nabla \psi_{\mathrm{dcon}}\right) \\[4pt]
&= \left(\frac{\chi^\prime_{\mathrm{dcon}}}{2\pi}\right)^2
|\nabla \psi_{\mathrm{dcon}}|^2 \sigma S
+ \sigma \mathbf{j}\!\cdot\!\mathbf{B}
+ p'\,(\boldsymbol{\kappa}\!\cdot\!\nabla \psi_{\mathrm{dcon}}) \chi^\prime / \pi.
\end{align}

Therefore, the final form of \(K\) in the DCON coordinate system is
\begin{equation}
\boxed{
K
= \left(\frac{\chi^\prime_{\mathrm{dcon}}}{2\pi}\right)^2
|\nabla \psi_{\mathrm{dcon}}|^2\,\sigma\,S
+ \sigma \mathbf{j}\!\cdot\!\mathbf{B}
+ p'\,(\boldsymbol{\kappa}\!\cdot\!\nabla \psi_{\mathrm{dcon}}) \chi^\prime / \pi.
}
\end{equation}


\subsection{Normalization of S between Chance and DCON coordinates}

In Chance's formulation, the magnetic shear scalar \(S\) is defined as
\begin{equation}
\color{blue} 
S
= -\,\frac{\mathbf{B}\times\nabla\psi}{|\nabla\psi|^2}
\cdot
\left(
\nabla\times
\frac{\mathbf{B}\times\nabla\psi}{|\nabla\psi|^2}
\right).
\end{equation}

Using the DCON normalization of the flux coordinate,
\begin{equation}
\color{blue} 
\psi
\color{black} 
= \frac{\chi^\prime_{\mathrm{dcon}}}{2\pi}\,
\psi_{\mathrm{dcon}},
\qquad
\color{blue} \nabla \psi
\color{black} 
= \frac{\chi^\prime_{\mathrm{dcon}}}{2\pi}\, \nabla \psi_{\mathrm{dcon}},
\end{equation}
we can rewrite each factor in terms of DCON variables.

Hence, the normalized cross term becomes
\begin{equation}
\frac{\mathbf{B}\times\nabla\psi}{|\nabla\psi|^2}
= \frac{2\pi}{\chi^\prime_{\mathrm{dcon}}}\,
\frac{\mathbf{B}\times\nabla\psi_{\mathrm{dcon}}}{|\nabla\psi_{\mathrm{dcon}}|^2}.
\end{equation}

Substituting this back into the original definition of \(S\), we find
\begin{align}
S
&= -\left(\frac{2\pi}{\chi^\prime_{\mathrm{dcon}}}\right)^2
\frac{\mathbf{B}\times\nabla\psi_{\mathrm{dcon}}}{|\nabla\psi_{\mathrm{dcon}}|^2}
\cdot
\left[
\nabla\times
\frac{\mathbf{B}\times\nabla\psi_{\mathrm{dcon}}}{|\nabla\psi_{\mathrm{dcon}}|^2}
\right]\\
&=-(2\pi)^2
\frac{\mathbf{B}\times\nabla\psi_{\mathrm{dcon}}\chi^\prime_{dcon}}{|\nabla\psi_{\mathrm{dcon}}*\chi^\prime_{dcon}|^2}
\cdot
\left[
\nabla\times
\frac{\mathbf{B}\times\nabla\psi_{\mathrm{dcon}}\chi^\prime_{dcon}}{|\nabla\psi_{\mathrm{dcon}}\chi^\prime_{dcon}|^2}
\right]\\
&= -\,(2\pi)^2 \frac{\mathbf{B}\times\nabla\chi_{dcon}}{|\nabla\chi_{dcon}|^2}
\cdot
\left(
\nabla\times
\frac{\mathbf{B}\times\nabla\chi_{dcon}}{|\nabla\chi_{dcon}|^2}
\right)
\end{align}

\begin{align}
\boldsymbol{\Lambda}
&\equiv \frac{\mathbf{B} \times \nabla \chi_{dcon}}{|\nabla \chi_{dcon}|^2}
= \frac{(\nabla \chi_{dcon} \times \nabla \alpha) \times \nabla \chi_{dcon}}{|\nabla \chi_{dcon}|^2} \\[4pt]
&= \nabla \alpha
- \frac{(\nabla \chi_{dcon} \cdot \nabla \alpha)}{|\nabla \chi_{dcon}|^2}\, \nabla \chi_{dcon}.
\end{align}

The curl becomes ($\alpha\, \text{represents}\, \alpha_{dcon}$ below this point)
\begin{align}
\nabla \times \boldsymbol{\Lambda}
&= -\,\nabla \times
\left[
\frac{(\nabla \chi_{dcon} \cdot \nabla \alpha)}{|\nabla \chi_{dcon}|^2}\, \nabla \chi_{dcon}
\right] \\[4pt]
&= -\,\nabla \left(\frac{\nabla \chi_{dcon} \cdot \nabla \alpha}{|\nabla \chi_{dcon}|^2}\right)
\times \nabla \chi_{dcon},
\end{align}

since $\nabla \times \nabla = 0$. The magnetic shear:

\begin{align}
S
&= - (2\pi)^2\, \boldsymbol{\Lambda} \cdot (\nabla \times \boldsymbol{\Lambda}) \\[4pt]
&= - (2\pi)^2
\left[
\nabla \alpha
- \frac{(\nabla \chi_{dcon} \cdot \nabla \alpha)}{|\nabla \chi_{dcon}|^2}\, \nabla \chi_{dcon}
\right]
\cdot
\left[
\nabla \left(\frac{\nabla \chi_{dcon} \cdot \nabla \alpha}{|\nabla \chi_{dcon}|^2}\right)
\times \nabla \chi_{dcon}
\right] \\[6pt]
&= (2\pi)^2\, \nabla \alpha \cdot
\left[
\nabla \left(\frac{\nabla \chi_{dcon} \cdot \nabla \alpha}{|\nabla \chi_{dcon}|^2}\right)
\times \nabla \chi_{dcon}
\right] \\[6pt]
&= (2\pi)^2\, (\nabla \chi_{dcon} \times \nabla \alpha) \cdot
\nabla \left(\frac{\nabla \chi_{dcon} \cdot \nabla \alpha}{|\nabla \chi_{dcon}|^2}\right) \\[6pt]
&= (2\pi)^2\, \mathbf{B} \cdot
\nabla \left(\frac{\nabla \chi_{dcon} \cdot \nabla \alpha}{|\nabla \chi_{dcon}|^2}\right)\\
&= \frac{(2\pi)^2}{\chi\prime_{dcon}}\, \mathbf{B} \cdot \nabla (\frac{\nabla \psi \cdot \nabla \alpha}{\nabla\psi \cdot \nabla\psi})
\end{align}


The parallel gradient operator is:
\begin{equation}
\mathbf{B} \cdot \nabla
= B^{\theta} \frac{\partial}{\partial \theta}
+ B^{\zeta} \frac{\partial}{\partial \zeta}
= \frac{\chi'}{\mathcal{J}}
\left(
\frac{\partial}{\partial \theta}
+ q \frac{\partial}{\partial \zeta}
\right).
\end{equation}

Also note
\begin{align}
\nabla \psi \cdot \nabla \alpha
&= \nabla \psi \cdot \nabla \big(q(\psi)\theta - \zeta\big)
= q'(\nabla \psi \cdot \nabla \psi)\theta
+ \nabla \psi \cdot \big(q \nabla \theta - \nabla \zeta\big).
\end{align}

Since $\partial / \partial \zeta = 0$ for a tokamak, one obtains:
\begin{align}
S
&= \frac{(2\pi)^2}{\mathcal{J}}
\frac{\partial}{\partial \theta}
\left[
q' \theta
+ \frac{q(\nabla \psi \cdot \nabla \theta)
- (\nabla \psi \cdot \nabla \zeta)}
{(\nabla \psi \cdot \nabla \psi)}
\right] \\[6pt]
&= \frac{(2\pi)^2}{\mathcal{J}}
\left[
q'
+ \frac{\partial}{\partial \theta}
\left(
\frac{q(\nabla \psi \cdot \nabla \theta)
- (\nabla \psi \cdot \nabla \zeta)}
{(\nabla \psi \cdot \nabla \psi)}
\right)
\right].
\end{align}

\subsection{$\xi^{\psi}$ in DCON berstein caculation}

\begin{align}
\boxed{
\boldsymbol{\xi}_\perp \cdot \hat{n}
= \frac{\boldsymbol{\xi} \cdot \nabla \psi}{|\nabla \psi|}
= \frac{\xi^\psi}{|\nabla \psi|},
}
\qquad
\boxed{
(\boldsymbol{\xi}_\perp \cdot \hat{n})^2
= \frac{(\xi^\psi)^2}{|\nabla \psi|^2}.
}
\end{align}

\begin{equation}
\boxed{
\boldsymbol{\xi}_\perp \cdot \hat{n}
= \frac{1}{|\nabla \psi|}
\sum_{m,n} v_{m,n}(\psi)\, e^{2\pi i(m\theta - n\zeta)}.
}
\end{equation}

\subsection{$\boldsymbol{\kappa}$ in DCON berstein caculation}


\section{Reconstruction of energy principle variables in DCON}
\subsection{Reconstruction v1}

\begin{align}
\xi^\psi = \sum_{m}^{} \sum_{n}^{} v_{m,n}(\psi) \cdot \exp 2\pi i(m\theta -n\zeta) \\
\xi^s = \sum_{m}^{} \sum_{n}^{} u_{m,n}(\psi) \cdot \exp 2\pi i(m\theta -n\zeta)
\end{align}

\begin{align}
\vec{j} &= \mathcal{J} (j^{\psi} \nabla \theta \times \nabla \zeta + j^{\theta} \nabla \zeta \times \nabla \psi + j^{\zeta} \nabla \psi \times \nabla \theta) \\
&= \mathcal{J} ( j^{\psi} \vec{v_1} + j^{\theta} \vec{v_2} + j^{\zeta} \vec{v_3})\\
&=\mathcal{J} (j^{\theta} \nabla \zeta \times \nabla \psi + j^{\zeta} \nabla \psi \times \nabla \theta)
\end{align}

\begin{equation}
\vec{Q} = \mathcal{J} (Q^{\psi} \nabla \theta \times \nabla \zeta + Q^{\theta} \nabla \zeta \times \nabla \psi + Q^{\zeta} \nabla \psi \times \nabla \theta) = \mathcal{J} ( Q^{\psi} \vec{v_1} + Q^{\theta} \vec{v_2} + Q^{\zeta} \vec{v_3})
\end{equation}

\begin{equation}
Q^{\psi} = \frac{\chi^\prime}{\mathcal{J}} (\frac{\partial}{\partial \theta} + q \frac{\partial}{\partial \zeta} )\xi^\psi = \frac{\chi^\prime}{\mathcal{J}} \sum_{m}^{} \sum_{n}^{} v_{m,n}(\psi) \cdot \exp 2 \pi i(m\theta -n\zeta)\cdot 2 \pi i (m-nq) 
\end{equation}

\begin{align}
Q^{\theta} &= - \frac{1}{\mathcal{J}} \frac{\partial}{\partial \psi}(\chi ^ \prime \xi ^\psi) + \frac{1}{\mathcal{J}} \frac{\partial \xi^s}{\partial \zeta} \\
&= - \frac{1}{\mathcal{J}} \sum_{m}^{} \sum_{n}^{} \left(\chi ^ \prime \frac{\partial}{\partial \psi}(v_{m,n}(\psi)) + 2\pi i n u_{m,n}(\psi) \right) \cdot \exp 2\pi i(m\theta -n\zeta)
\end{align}

\begin{align}
Q^{\zeta}
&= - \frac{1}{\mathcal{J}} \frac{\partial}{\partial \psi} (\chi^\prime q \xi^\psi)
   -\frac{1}{\mathcal{J}} \frac{\partial \xi^s}{\partial \theta} \\
&= - \frac{1}{\mathcal{J}} \sum_{m} \sum_{n}
   \left( \chi ^ \prime q \frac{\partial v_{m,n}(\psi)}{\partial \psi}
   + \chi^\prime q^\prime v_{m,n}(\psi)
   + 2\pi i m u_{m,n}(\psi) \right)
   \exp\!\bigl(2 \pi i(m\theta - n\zeta)\bigr).
\end{align}

\begin{equation}
\Xi_{\psi} 
= \texttt{u(:,:,1)}  \quad \text{(Fortran, DCON)}
= \texttt{u(:,:,1)}  \quad \text{(Fortran, MATCH)}.
\end{equation}

\begin{equation}
\Xi_{s}
= \texttt{ud(:,:,2)}  \quad \text{(Fortran, DCON)}
= \texttt{u(:,:,4)}  \quad \text{(Fortran, MATCH)}.
\end{equation}

\begin{equation}
\Xi_{\psi}^\prime
= \texttt{ud(:,:,1)}
= \texttt{du(:,:,1)} \quad \text{(Fortran, DCON)}
= \texttt{u(:,:,3)}  \quad \text{(Fortran, MATCH)}.
\end{equation}

And if you see the 331 line in ideal.f of MATCH, it constructs all the eigenvectors in variable v.
\begin{equation}
v_{m,n}(\psi)
= v(:,1,mstep)
\end{equation}

\begin{equation}
v_{m,n}^\prime(\psi)
= v(:,3,mstep)
\end{equation}

\begin{equation}
u_{m,n}(\psi)
= v(:,4,mstep)
\end{equation}

\begin{equation}
    \mu_0 \vec{j} = \mathcal{J} (\mu_0 j^\theta \vec{v}_2 + \mu_0 j^\zeta \vec{v}_3)
\end{equation}
\begin{equation}
    \mu_0 \vec{j} \times \nabla \psi = \mathcal{J} \mu_0 j^\theta (\vec{v}_2 \times \nabla \psi) + \mathcal{J} \mu_0 j^\zeta (\vec{v}_3 \times \nabla \psi)
\end{equation}

\begin{equation}
    j^\theta = j \cdot \nabla \theta = - 2 \pi f^\prime / \mathcal{J}
\end{equation}

\begin{equation}
    j^\zeta = j \cdot \nabla \zeta = - \mu_0 p^\prime / \chi^\prime - 2 \pi f^\prime q / \mathcal{J}
\end{equation}
Now we can reconstruct the cross term in the energy principle:

\begin{align}
\vec{C} &\equiv \vec{Q} + (\vec{\xi} \cdot \hat{n} ) (\mu_0 \vec{j} \times \hat{n}) \\
&=C^\psi \vec{e}_1 + C^\theta \vec{e}_2 + C^\zeta \vec{e}_3 = \mathcal{J} ( C^\psi \vec{v}_1 + C^\theta \vec{v}_2 + C^\zeta \vec{v}_3 )\\
&=C_1 \nabla\psi + C_2 \nabla\theta + C_3 \nabla\zeta
\end{align}

\begin{align}
C^\psi &= \vec{C} \cdot \nabla\psi = Q^\psi \\
C^\theta &= \vec{C} \cdot \nabla\theta = Q^\theta + \frac{\xi^\psi}{|\nabla\psi|^2} \mu_0 \vec{j} \cdot (\nabla\psi \times \nabla\theta) \\
C^\zeta &= \vec{C} \cdot \nabla\zeta = Q^\zeta + \frac{\xi^\psi}{|\nabla\psi|^2} \mu_0 \vec{j} \cdot (\nabla\psi \times \nabla\zeta)
\end{align}

\begin{align}
C_1 &= Q_1 + \frac{\mathcal{J}\xi^\psi}{|\nabla\psi|^2} \left[ \mu_0 j^\theta (\nabla\psi \cdot \nabla\zeta) - \mu_0 j^\zeta (\nabla\psi \cdot \nabla\theta) \right] \\
C_2 &= Q_2 + \mathcal{J}\xi^\psi \mu_0 j^\zeta \\
C_3 &= Q_3 - \mathcal{J}\xi^\psi \mu_0 j^\theta
\end{align}

\begin{align}
\delta W_{term1} &= \frac{1}{2} \int \bigl(\frac{|\vec{C}|^2}{\mu_0} \bigr) dx^3 \\
&= \frac{1}{2} \int d\psi d\theta d\zeta \mathcal{J} \Bigl[ \frac{|\vec{C}|^2}{\mu_0}\Bigr]
\end{align}

\begin{align}
g_{11} &\equiv \mathcal{J}^2\,|\nabla\theta\times\nabla\zeta|^2, \\
g_{22} &\equiv \mathcal{J}^2\,|\nabla\zeta\times\nabla\psi|^2, \\
g_{33} &\equiv \mathcal{J}^2\,|\nabla\psi\times\nabla\theta|^2, \\
g_{12} &\equiv \mathcal{J}^2\,(\nabla\theta\times\nabla\zeta)\cdot(\nabla\zeta\times\nabla\psi), \\
g_{23} &\equiv \mathcal{J}^2\,(\nabla\zeta\times\nabla\psi)\cdot(\nabla\psi\times\nabla\theta), \\
g_{31} &\equiv \mathcal{J}^2\,(\nabla\psi\times\nabla\theta)\cdot(\nabla\theta\times\nabla\zeta),
\end{align}
so that
\begin{equation}
\boxed{\;\mathcal{J}\,(\vec v_i\cdot\vec v_j)=\frac{g_{ij}}{\mathcal{J}}\;}
\end{equation}

\begin{equation}
\boxed{
\left\langle F \right\rangle_{mm'}
\equiv \int_{0}^{1} d\theta\;
\mathcal{J}(\psi,\theta)\,
\exp\!\bigl(2\pi i(m'-m)\theta\bigr)\,
F(\psi,\theta)
}
\end{equation}

\begin{equation}
\boxed{
(G_{ij})_{mm'} \equiv \left\langle \frac{g_{ij}}{\mathcal{J}} \right\rangle_{mm'}
}
\end{equation}

\begin{align}
\delta W_{\mathrm{term1}}
&= \frac{1}{2} \int \bigl(\frac{|\vec{C}|^2}{\mu_0} \bigr) dx^3 \\
&= \frac{1}{2\mu_0}\int d\psi\,d\theta\,d\zeta\;
\mathcal{J}\,|\vec C|^2
= \frac{1}{2\mu_0}\int d\psi\,d\theta\,d\zeta\;\mathcal{J} \sum_{i,j=1}^{3} C^{i} C^{j*} g_{ij}\\
&=\frac{1}{2\mu_0} \int d\psi d\theta\,d\zeta\ \left[ \mathcal{J} \sum_{i,j}\ \sum_{m,n}\ \sum_{m'n'}\ C_{m,n}^{i} C_{m',n'}^{j*} g_{ij} e^{2\pi i ((m-m')\theta -(n-n') \zeta)}\right] \\
\end{align}

Since $n$ is fixed in DCON so,
\begin{align}
    \delta W_{n=fixed}&=\frac{1}{2\mu_0} \int d\psi\ \left[ \int d\theta \, \sum_{m, m'} e^{2\pi i(m-m')\theta} \mathcal{J}(\psi, \theta) \sum_{i,j} C_{m,n}^{i}(\psi, \theta) C_{m',n}^{j*}(\psi, \theta) g_{ij}(\psi, \theta) \right] \\
    &=\frac{1}{2\mu_0} \int d\psi \int_0^1 d\theta \; \mathcal{J}(\psi, \theta) \sum_{i,j} \left( \sum_m C_{m,n}^i e^{2\pi i m \theta} \right) \left( \sum_{m'} C_{m',n}^{j*} e^{-2\pi i m' \theta} \right) g_{ij}(\psi, \theta)
\end{align}

The second term is $\nabla \cdot \vec{\xi} = 0$, it vanishes:
\begin{equation}
\boxed{
\delta W_{\mathrm{term2}} 
= \frac{1}{2} \int \gamma p |\vec{\nabla} \cdot \vec{\xi}|^2 dx^3 = \frac{1}{2} \int d\psi d\theta d\zeta \mathcal{J} \Bigl[ \gamma p |\vec{\nabla} \cdot \vec{\xi}|^2 \Bigr] = 0.
}
\end{equation}

For the last term in $\delta W$,
\begin{align}
   \delta W_{\mathrm{term3}} &= \frac{1}{2} \int -K(\vec{\xi} \cdot \hat{n})^2 dx^3 \\
   &= \frac{1}{2} \int d\psi d\theta d\zeta \mathcal{J} \Bigl[ -K(\vec{\xi} \cdot \hat{n})^2 \Bigr] \\
   &= \frac{1}{2} \int d\psi d\theta d\zeta \mathcal{J} \Bigl[ -K(\frac{\vec{\xi}^{\psi}}{|\nabla\psi|})^2 \Bigr]
\end{align}

% --- Term 3 as quadratic form ---
\begin{equation}
\delta W_{\mathrm{term3}}
= -\frac{1}{2}\int d\psi\,d\theta\,d\zeta\;
\mathcal{J}\,\frac{K}{|\nabla\psi|^2}\;
\xi^\psi\,\xi^{\psi *}.
\end{equation}

\begin{equation}
(K_{mm'}) \equiv \left\langle \frac{K}{|\nabla\psi|^2}\right\rangle_{mm'}
= \int_0^1 d\theta\;\mathcal{J}(\psi,\theta)\,
e^{2\pi i(m'-m)\theta}\,
\frac{K}{|\nabla\psi|^2}.
\end{equation}

\begin{equation}
\boxed{
\delta W_{\mathrm{term3}}
= -\frac{1}{2}\int d\psi \sum_{m,m'}
v^*_{m,n}(\psi)\;
(K_{mm'}(\psi))\;
v_{m',n}(\psi).
}
\end{equation}

\subsection{Reconstruction $\boldsymbol{v2}$ : Matrix representation of $\delta W$}


\begin{equation}
\begin{aligned}
2\delta W 
&= \int \big[ \frac{|\vec{C}|^2}{\mu_0}- K(\vec{\xi} \cdot \hat{n})^2 +\gamma p |\vec{\nabla} \cdot \vec{\xi}|^2 \big] dx^3\\
&= \int d\psi d\theta d\zeta \mathcal{J} \Big\{ [  \nabla \xi_s \times \nabla \psi
+ (\nabla \theta \times \nabla \zeta)\,\mathcal{J}\,\mathbf{B}\cdot\nabla \xi_\psi
- \mathbf{B}\,\xi_\psi' \\
&-
\left(
    \chi'' \nabla \zeta \times \nabla \psi
    + (q\chi')' \nabla \psi \times \nabla \theta
\right)\xi_\psi \\
&+\frac{\xi_\psi}{|\nabla \psi|^2} [
j^\theta
\big(
\nabla\psi(\nabla\zeta\cdot\nabla\psi)
-
\nabla\zeta|\nabla\psi|^2
\big)+j^\zeta
\big(
\nabla\theta|\nabla\psi|^2
-
\nabla\psi(\nabla\theta\cdot\nabla\psi)
\big)]]^2 \\
&-K(\vec{\xi} \cdot \hat{n})^2 +\gamma p |\vec{\nabla} \cdot \vec{\xi}|^2\}
\end{aligned}
\end{equation}

\subsubsection{Matrix representation of $\delta W_{\mathrm{term1}}$}
First we focus on the first term in $\delta W$. If we change each terms in $\vec{C}$ as follows:
\begin{equation}
\begin{aligned}
&\mathbf X \equiv \nabla \xi_s \times \nabla \psi,\\
&\mathbf Y \equiv (\nabla \theta \times \nabla \zeta)\, \mathcal J \, \mathbf B \cdot \nabla \xi_\psi,\\
&\mathbf Z \equiv -\mathbf B\, \xi_\psi',\\
&\mathbf U \equiv -\left(
\chi'' \nabla \zeta \times \nabla \psi
+
(q\chi')'
\nabla \psi \times \nabla \theta
\right)\xi_\psi,\\
&\mathbf V \equiv
\frac{\xi_\psi}{|\nabla \psi|^2}
\, \mu_0 \mathbf j \times \nabla \psi .    
\end{aligned}
\end{equation}

\begin{equation}
\begin{aligned}
2\delta W_{\mathrm{term1}} & = \int d\psi d\theta d\zeta \mathcal{J} [\frac{|\vec{C}|^2}{\mu_0}]\\
&= \int d\psi d\theta d\zeta \mathcal{J}|\mathbf C|^2\\
&= \int d\psi d\theta d\zeta \mathcal{J} \{ |\mathbf Q|^2+
2\,\Re\!\left[
(\mathbf X+\mathbf Y+\mathbf Z+\mathbf U)\cdot\mathbf V^*
\right]+|\mathbf V|^2\} 
\end{aligned}
\end{equation}

The detailed derivation of the term already in \cite{10.1063/1.4958328} willbe skipped here.
The first term $|\mathbf X|^2$ can be written with matrix form as $\Xi_s^\dagger \mathbf A \Xi_s$ where $\Xi_s$ is a vector of the coefficients of the poloidal harmonics in $\xi_s$, and $\mathbf A$ is a matrix given by:
\begin{equation}
\mathbf A=(2\pi)^2\left[n^2 \mathbf G_{22}+n \mathbf G_{23} \mathbf M+n \mathbf M \mathbf G_{23}+\mathbf M \mathbf G_{33} \mathbf M\right].
\end{equation}

The second term $\mathbf X \cdot \mathbf Z + \mathbf Z \cdot \mathbf X$ can be written in matrix form as $\Xi_s^\dagger \mathbf B \Xi_\psi^\prime + \Xi_\psi'^\dagger \mathbf B^\dagger \Xi_s$, where $\mathbf B$ is a matrix given by:
\begin{equation}
    \mathbf B = -2\pi i\chi' [n(\mathbf G_{22} + q\mathbf G_{23}) + \mathbf M(\mathbf G_{23} + q\mathbf G_{33})]
\end{equation}

The term $\mathbf X \cdot \mathbf Y + \mathbf Y \cdot \mathbf X$ as $\Xi_s^\dagger \mathbf C_1 \Xi_\psi + \Xi_\psi^\dagger \mathbf C_1^\dagger \Xi_s$, where $\mathbf C_1$ is a matrix given by:
\begin{equation}
    \mathbf C_1 = -(2\pi)^2 \chi^\prime (\mathbf M \mathbf G_{31} + n \mathbf G_{12})\mathbf Q
\end{equation}

The term $\mathbf X \cdot \mathbf U + \mathbf U \cdot \mathbf X$ as $\Xi_s^\dagger \mathbf C_2 \Xi_\psi + \Xi_\psi^\dagger \mathbf C_2^\dagger \Xi_s$, where $\mathbf C_2$ is a matrix given by:
\begin{equation}
    \mathbf C_2 = -2\pi [\chi'' (n \mathbf G_{22} + \mathbf M \mathbf G_{23}) + (q\chi')' (n \mathbf G_{23} + \mathbf M \mathbf G_{33})] 
\end{equation}

The term $|\mathbf Y|^2$ as $\Xi_\psi^\dagger \mathbf H_1 \Xi_\psi$, where $\mathbf H_1$ is a matrix given by:
\begin{equation}
    \mathbf H_1 = (2\pi\chi^{\prime} )^2 \mathbf Q\mathbf G_{11} \mathbf Q
\end{equation}

The term $\mathbf Y \cdot \mathbf Z + \mathbf Z \cdot \mathbf Y$ as $\Xi_\psi^{\prime\dagger }\mathbf E_1 \Xi_\psi + \Xi_\psi^{\dagger} \mathbf E_1^\dagger \Xi_\psi^\prime$, where $\mathbf E_1$ is a matrix given by:
\begin{equation}
    \mathbf E_1 = -2\pi i \chi^{\prime 2} (\mathbf G_{12}+q \mathbf G_{31})\mathbf Q
\end{equation}

The term $\mathbf Y \cdot \mathbf U + \mathbf U \cdot \mathbf Y$ as $\Xi_\psi^\dagger \mathbf H_2 \Xi_\psi$, where $\mathbf H_2$ is a matrix given by:
\begin{equation}
    \mathbf H_2 = 2\pi i \chi'
\left[
\chi''\left(\mathbf M \mathbf G_{12} -\mathbf G_{12} \mathbf M\right)
+
(q\chi')'\left(\mathbf M \mathbf G_{31} -\mathbf G_{31} \mathbf M\right)
\right]
\end{equation}

The term $|\mathbf Z|^2$ as $\Xi_\psi^{\prime \dagger} \mathbf D \Xi_\psi^\prime$, where $\mathbf D$ is a matrix given by:
\begin{equation}
    \mathbf D = \chi'^2
\left[
(\mathbf{G}_{22} + q\,\mathbf{G}_{23})
+
q\,(\mathbf{G}_{23} + q\,\mathbf{G}_{33})
\right]
\end{equation}

The term $\mathbf Z \cdot \mathbf U + \mathbf U \cdot \mathbf Z$ as $\Xi_\psi^{\prime\dagger}\mathbf E_2 \Xi_\psi + \Xi_\psi^\dagger \mathbf E_2^\dagger \Xi_\psi^\prime$, where $\mathbf E_2$ is a matrix given by:
\begin{equation}
\mathbf E_2 = \chi'
\left[
\chi''
\left(
\mathbf{G}_{22}
+
q\,\mathbf{G}_{23}
\right)
+
(q\chi')'
\left(
\mathbf{G}_{23}
+
q\,\mathbf{G}_{33}
\right)
\right]
\end{equation}

Finally, the term $|\mathbf U|^2$ as $\Xi_\psi^\dagger \mathbf H_3 \Xi_\psi$, where $\mathbf H_3$ is a matrix given by:
\begin{equation}
\mathbf{H}_3
=
\chi''
\left[
\chi''\,\mathbf{G}_{22}
+
(q\chi')'\,\mathbf{G}_{23}
\right]
+
(q\chi')'
\left[
\chi''\,\mathbf{G}_{23}
+
(q\chi')'\,\mathbf{G}_{33}
\right]
\end{equation}

The new terms are the cross terms between $\mathbf V$ and the other four terms. The term $\mathbf X \cdot \mathbf V + \mathbf V \cdot \mathbf X$ can be written in matrix form as $\Xi_s^\dagger \mathbf C_3 \Xi_\psi + \Xi_\psi^\dagger \mathbf C_3^\dagger \Xi_s$, where $\mathbf C_3$ is a matrix given by:
\begin{equation}
(\nabla \xi_s \times \nabla \psi) \cdot (\frac{\xi_\psi}{|\nabla \psi|^2} \mu_0 \mathbf j \times \nabla \psi) = (\nabla \xi_s \cdot \mu_0 \mathbf j) - \cancel{(\nabla \xi_s \cdot \nabla \psi)(\mu_0 \mathbf j \cdot \nabla \psi)}
\end{equation}

\begin{align}
    \mathbf{X}^* \cdot \mathbf{V}&=
\left[
-2\pi i\, \xi_s^*
\left(
m\, \nabla\theta \times \nabla\psi
-
n\, \nabla\zeta \times \nabla\psi
\right)
\right]
\cdot
\left[
\frac{\xi_\psi}{|\nabla\psi|^2}
\,\mu_0
\left(
\mathbf{j} \times \nabla\psi
\right)
\right]\\
&=
(-2\pi i)\,\mu_0\,
\xi_s^*\,
\left(
m\, j_\theta
-
n\, j_\zeta
\right)
\xi_\psi\,
\end{align}
This term is same as the term $\xi_\psi \mathbf J \cdot \nabla \psi - \xi_s \mathbf J \cdot \nabla \psi$ in original form. 
And the matrix form can be written as $\Xi_s^\dagger \mathbf C_3 \Xi_\psi + \Xi_\psi^\dagger \mathbf C_3^\dagger \Xi_s$, where $\mathbf C_3$ is a matrix given by:
\begin{equation}
\mathbf C_3 = 2\pi i (2\pi f^\prime \mathbf Q -n \frac{p^\prime}{\chi^\prime} \mathbf J)
\end{equation}

The term $Z \cdot V + V \cdot Z$ can be written in matrix form as $\Xi_\psi^\dagger \mathbf E_3 \Xi_\psi^\prime + \Xi_\psi^{\dagger\prime} \mathbf E_3^\dagger \Xi_\psi + \Xi_\psi^\dagger \mathbf H_4 \Xi_\psi$, where $\mathbf E_3$ is a matrix given by:

\begin{equation}
\mathbf Y^*
=
(\nabla \theta \times \nabla \zeta)\,
(-2\pi i)\chi'
\sum_m
(m - nq)\,
v_m^*\,
e^{-2\pi i (m\theta - n\zeta)} .
\end{equation}

\begin{equation}
\mathbf V
=
\sum_{m'}
v_{m'}\,
e^{2\pi i (m'\theta - n\zeta)}
\,
\frac{\mu_0 \mathbf j \times \nabla \psi}
{|\nabla \psi|^2} .
\end{equation}

\begin{equation}
\begin{aligned}
\mathbf Y^* \cdot \mathbf V
&=
(-2\pi i)\chi'
\sum_{m,m'}
v_m^*
(m - nq)
v_{m'}
e^{2\pi i (m' - m)\theta}
\;
\frac{
(\nabla \theta \times \nabla \zeta)
\cdot
(\mu_0 \mathbf j \times \nabla \psi)
}{
|\nabla \psi|^2
}\\
&=
(-2\pi i)\chi'
\sum_{m,m'}
v_m^*
(m - nq)
v_{m'}
e^{2\pi i (m' - m)\theta}
\;
\mu_0 \,
\frac{
j^{\theta} \left( \nabla \zeta \cdot \nabla \psi \right)
-
j^{\zeta} \left( \nabla \theta \cdot \nabla \psi \right)
}{
|\nabla \psi|^2
}.
\end{aligned}
\end{equation}


\section{Plasma Magnetic Potential on the Surface}

The plasma magnetic potential $\Phi$ on the boundary surface is computed using covariant magnetic field components obtained from the perturbed equilibrium solution. The calculation involves four components that represent different aspects of the magnetic potential response.

\subsection{Equilibrium Data}

The equilibrium spline $s_q$ contains the following flux functions indexed by their position in the spline:
\begin{align}
s_q(\psi) &= \begin{cases}
s_q(\psi)_1 &: \psi_{\text{fac}} \; (\text{flux normalization}) \\
s_q(\psi)_2 &: F(\psi) = R B_\phi \; (\text{poloidal current}) \\
s_q(\psi)_3 &: \mu_0 p(\psi) \; (\text{pressure term}) \\
s_q(\psi)_4 &: q(\psi) \; (\text{safety factor}) \\
s_q(\psi)_5 &: \rho(\psi) \; (\text{auxiliary quantity})
\end{cases}
\end{align}

\subsection{Covariant Field Components Transformation}

From the contravariant field components in mode space $(B_{w,p}^{(m)}, B_{m,t}^{(m)}, B_{m,z}^{(m)})$, the covariant components in physical space are obtained via inverse Fourier transform and metric tensor contraction:
\begin{align}
B_{v,\theta} &= \mathcal{F}^{-1}\{b_{v,t}^{(m)}\} \quad \text{(covariant theta component)}\\
B_{v,\zeta} &= \mathcal{F}^{-1}\{b_{v,z}^{(m)}\} \quad \text{(covariant zeta component)}\\
\Xi_\psi &= \mathcal{F}^{-1}\{x_{w,p}^{(m)}\} \quad \text{(poloidal flux perturbation)}
\end{align}

where $\mathcal{F}^{-1}$ denotes the inverse discrete Fourier transform in the poloidal angle direction.

\subsection{Magnetic Potential Components}

The four components of the plasma magnetic potential on the surface are computed as follows:

\begin{equation}
\begin{aligned}
\Phi_3(\theta) &= -\frac{B_{v,\zeta}(\theta)}{2\pi i n}
\\
\Phi_1(\theta) &= \Phi_3(\theta) - \frac{\frac{dF}{d\psi} \cdot \Xi_\psi(\theta)}{2\pi i n \cdot \mathcal{J}(\theta)}
\\
\Phi_4(\theta) &= B_{v,\theta}(\theta)
\\
\Phi_2(\theta) &= \Phi_4(\theta) - \left(
\frac{dF}{d\psi} q(\psi) + \mathcal{J}(\theta) \frac{d(\mu_0 p)}{d\psi} \chi_1^{-1}
\right) \frac{\Xi_\psi(\theta)}{\mathcal{J}(\theta)}
\end{aligned}
\end{equation}

Here:
\begin{itemize}
\item $\Phi_3$ represents the potential associated with toroidal field perturbation ($\zeta$ component)
\item $\Phi_1$ includes the coupling between flux perturbation and poloidal current derivative
\item $\Phi_4$ is the direct poloidal field response  
\item $\Phi_2$ combines poloidal field response with pressure and current effects
\item $\mathcal{J}(\theta) = \partial(\mathbf{r}, \mathbf{z}, \phi)/\partial(\psi, \theta, \zeta)$ (Jacobian of coordinate transformation)
\item $\chi_1 = \psi_0 \cdot 2\pi$ (flux scaling factor)
\item $n$ is the toroidal mode number
\end{itemize}

\subsection{Physical Interpretation}

Each potential component couples different physical mechanisms:
\begin{enumerate}
\item \textbf{$\Phi_3$}: Direct response to toroidal field perturbation, normalized by mode number.

\item \textbf{$\Phi_1$}: Additional correction accounting for the interaction between the poloidal flux perturbation $\Xi_\psi$ and the poloidal current profile $(dF/d\psi)$. The denominator $\mathcal{J}$ accounts for local metric effects.

\item \textbf{$\Phi_4$}: The primary poloidal magnetic field response from the perturbed equilibrium.

\item \textbf{$\Phi_2$}: Complex response combining:
\begin{itemize}
\item Gradient of poloidal current: $(dF/d\psi) q$ 
\item Pressure gradient effect: $\mathcal{J} (d\mu_0 p/d\psi)\chi_1^{-1}$
\item Both weighted by normalized flux perturbation: $\Xi_\psi/\mathcal{J}$
\end{itemize}
\end{enumerate}


\bibliographystyle{plain}
\bibliography{references}

\end{document}